\section{定义}\label{section:1}
\SourcePage{16}
\subsection{范畴}\label{sec:1a}
\begin{definition*}
固定系数域 $\K$。我们假定所有范畴都是 $\K$ 上的线性小范畴。一个\emph{非含幺 $\Ai$ 范畴} $\A$ 由对象集 $\Ob\A$、对每对对象给出的分次向量空间 $\hom_{\A}(X_0,X_1)$，以及对每个阶数 $d\ge1$ 给出的复合映射组成：
\begin{equation}\label{eq:1.1}
\mu^d_{\A}:\hom_{\A}(X_{d-1},X_d)\otimes\cdots\otimes\hom_{\A}(X_0,X_1)\longrightarrow\hom_{\A}(X_0,X_d)[2-d].
\end{equation}
照例，$[k]$ 表示将向量空间的分次向下平移 $k\in\Z$。映射\eref{1.1}必须满足以下二次的 $\Ai$ 结合律方程：
\begin{equation}\label{eq:1.2}
\sum_{m,n}(-1)^{\sgnmark_n}\mu^{d-m+1}_{\A}\bigl(a_d,\ldots,a_{n+m+1},\mu^m_{\A}(a_{n+m},\ldots,a_{n+1}),a_n,\ldots,a_1\bigr)=0,
\end{equation}
其中对所有可能的项求和，即 $1\le m\le d$、$0\le n\le d-m$，符号由 $\sgnmark_n=|a_1|+\cdots+|a_n|-n$ 决定。同一和式还会出现在许多其他符号表达式中，因此从现在起，我们将它规定为通用记号。\znote{8}{记号 $\sgnmark_n$ 是纯粹的 Koszul 记账：把元素逐个移过前面的态射，每跨过一个次数为 $e$ 的态射便得因子 $(-1)^{e}$；但 $\mu^d$ 的输入已带 $[2-d]$ 的次数平移，实际出现的因而是 $|a_i|-1$ 而非 $|a_i|$，$n$ 项累加便得 $|a_1|+\cdots+|a_n|-n$。后文凡出现同一记号都可直接代入此式，不必重新推导。[译者按（上下文推断，无外部依据）：由\eref{1.1}的 $[2-d]$ 平移与 Koszul 符号规则相加即得。]}
\end{definition*}

\begin{definition*}
对于任意非含幺 $\Ai$ 范畴 $\A$，可以赋予它一个\emph{上同调范畴} $H(\A)$。其对象与 $\A$ 相同；态射空间为上同调群 $H(\hom_{\A}(X_0,X_1),\mu^1_{\A})$；复合则定义为
\begin{equation}\label{eq:1.3}
[a_2]\cdot[a_1]=(-1)^{|a_1|}[\mu^2_{\A}(a_2,a_1)].
\end{equation}
除了可能缺少恒等态射外，这是一个通常的线性分次范畴。\znote{8}{\textquotedblleft 非含幺\textquotedblright 只是说定义里不把恒等态射列为数据，并不禁止它存在。原书后面区分严格含幺（单位元在链层次上与全部 $\mu^d$ 相容）与上同调含幺（单位元只需在 $H$ 中起作用）等情形，并在讲完恒等态射之后约定（§2h）一律假定 c-含幺。阅读后文时务必先确认当前处在哪种假设下，否则\textquotedblleft 同构\textquotedblright、\textquotedblleft 拟等价\textquotedblright 的说法会失准。[来源：Seidel, Fukaya categories and Picard–Lefschetz theory, §2h（术语约定）及 §2 对恒等态射的讨论。]}还有一个变体，即只保留零次上同调的子范畴 $H^0(\A)\subset H(\A)$。同样地，如果一个非含幺 $\Ai$ 范畴 $\A$ 满足对所有 $d>2$ 都有 $\mu^d_{\A}=0$，那么它与一个非含幺\emph{微分分次范畴}（dg 范畴）是同一回事。更精确地，相关 dg 范畴的微分为 $\partial a=(-1)^{|a|}\mu^1_{\A}(a)$，乘法为 $a_2a_1=(-1)^{|a_1|}\mu^2_{\A}(a_2,a_1)$，与\eref{1.3}一致。
\end{definition*}

\begin{definition*}
非含幺 $\Ai$ 范畴的\emph{反范畴}定义如下：令 $\Ob\A^{\mathrm{opp}}=\Ob\A$、$\hom_{\A^{\mathrm{opp}}}(X_0,X_1)=\hom_{\A}(X_1,X_0)$，并将复合映射中各项的顺序反转：
\[
\mu^d_{\A^{\mathrm{opp}}}(a_d,\ldots,a_1)=(-1)^{\sgnmark_d}\mu^d_{\A}(a_1,\ldots,a_d).
\]
在上同调层次上，这给出通常的反范畴（不过不带 Koszul 符号），即 $H(\A^{\mathrm{opp}})=H(\A)^{\mathrm{opp}}$。\znote{8}{反范畴公式里的 $(-1)^{\sgnmark_d}$ 只是把被反转的输入顺序重新记账所需的 Koszul 符号，并非新的数学内容：到了上同调层次，$\mu^2_{\A^{\mathrm{opp}}}$ 恰给出反向复合，符号自动抵消。后文凡见 $\A^{\mathrm{opp}}$（例如 $\nufun(\A^{\mathrm{opp}},\Ch)$）都指这一约定，与取对偶无关。[译者按（上下文推断，无外部依据）：把\eref{1.2}用于 $\A^{\mathrm{opp}}$，再与 $H(\A)^{\mathrm{opp}}$ 的复合式逐项比对即得。]}
\end{definition*}

\begin{note}\label{st:1.1}
关于术语和约定，补充几点说明：(i) “非含幺”其实应当理解为“不要求含幺”。(ii) 假定范畴是小范畴，偶尔会造成形式上的麻烦，例如在定义扭曲复形时；我们把必要的集合论处理留给读者。(iii) 使用没有恒等态射的范畴时，需要小心：
\SourcePage{17}
例如，此时没有合理的对象同构概念；全忠实函子仍有意义，范畴的同构也有意义，但范畴的等价却没有意义，至少不能直接按通常方式定义。(iv) 我们的符号遵循“条形约定”，这是代数学家最常使用的约定；\eref{1.1}中递减的编号，试图调和态射复合的标准约定与圆周上标记点同样标准的循环次序，见\secref{9c}。
\end{note}

\begin{note}\label{st:1.2}
阶数 $d>2$ 的复合不是链映射。不过，它们仍以 Massey 积的形式，部分体现在上同调层次上。设有态射 $[a_k]\in\Hom_{H(\A)}(X_{k-1},X_k)$，$k=1,2,3$，并且两个乘积 $[a_3]\cdot[a_2]$ 与 $[a_2]\cdot[a_1]$ 都为零。选择链层次的代表元 $a_k$，再取 $h_k\in\hom_{\A}(X_{k-1},X_{k+1})$，$k=1,2$，使 $\mu^1_{\A}(h_k)=\mu^2_{\A}(a_{k+1},a_k)$。那么
\[
c=\mu^3_{\A}(a_3,a_2,a_1)-\mu^2_{\A}(h_2,a_1)-\mu^2_{\A}(a_3,h_1)\in\hom_{\A}(X_0,X_3)
\]
是次数为 $|a_1|+|a_2|+|a_3|-1$ 的上闭链，而类
\begin{equation}\label{eq:1.4}
[c]\in\frac{\Hom_{H(\A)}(X_0,X_3)}{[a_3]\cdot\Hom_{H(\A)}(X_0,X_2)+\Hom_{H(\A)}(X_1,X_3)\cdot[a_1]}
\end{equation}
只依赖于最初给定的 $[a_k]$。通常记作 $\langle[a_3],[a_2],[a_1]\rangle=(-1)^{|a_2|}[c]$。
\end{note}

\subsection{函子}\label{sec:1b}
\begin{definition*}
两个非含幺 $\Ai$ 范畴 $\A,\B$ 之间的\emph{非含幺 $\Ai$ 函子}，由映射 $\F:\Ob\A\to\Ob\B$ 以及对所有 $d\ge1$ 给出的多重线性映射组成：
\begin{equation}\label{eq:1.5}
\F^d:\hom_{\A}(X_{d-1},X_d)\otimes\cdots\otimes\hom_{\A}(X_0,X_1)\longrightarrow\hom_{\B}(\F X_0,\F X_d)[1-d],
\end{equation}
其中 $X_0,\ldots,X_d\in\Ob\A$，并满足以下多项式方程：
\begin{equation}\label{eq:1.6}
\begin{aligned}
&\sum_r\sum_{s_1,\ldots,s_r}\mu^r_{\B}\bigl(\F^{s_r}(a_d,\ldots,a_{d-s_r+1}),\ldots,\F^{s_1}(a_{s_1},\ldots,a_1)\bigr)\\
&\quad=\sum_{m,n}(-1)^{\sgnmark_n}\F^{d-m+1}\bigl(a_d,\ldots,a_{n+m+1},\mu^m_{\A}(a_{n+m},\ldots,a_{n+1}),a_n,\ldots,a_1\bigr).
\end{aligned}
\end{equation}
左侧对所有 $r\ge1$ 以及分拆 $s_1+\cdots+s_r=d$ 求和。
\end{definition*}

\begin{definition*}
它诱导一个通常的线性分次非含幺函子 $H(\F):H(\A)\to H(\B)$，在态射上的作用为 $[a]\mapsto[\F^1(a)]$。若 $H(\F)$ 是范畴同构，就称 $\F$ 为\emph{拟同构}；若 $H(\F)$ 全忠实，就称 $\F$ 为\emph{上同调全忠实}的。
\end{definition*}

非含幺 $\Ai$ 函子可按如下方式复合：
\begin{equation}\label{eq:1.7}
\begin{aligned}
(\G\circ\F)^d(a_d,\ldots,a_1)
=\sum_r\sum_{s_1,\ldots,s_r}\G^r\bigl(\F^{s_r}(a_d,\ldots,a_{d-s_r+1}),\ldots,\F^{s_1}(a_{s_1},\ldots,a_1)\bigr),
\end{aligned}
\end{equation}
\SourcePage{18}
其中求和范围与\eref{1.6}相同。复合严格满足结合律，而恒等函子——$\F X=X$、$\F^1=\id$（作为各态射空间的自同态），所有高阶项为零——是复合的单位元。

\subsection{形式微分同胚}\label{sec:1c}
\begin{definition*}
下面从一个稍有不同的角度看待某一特殊类的 $\Ai$ 函子。设有对象集 $\Ob\A$，并对每对对象给定分次向量空间 $\hom_{\A}(X_0,X_1)$，但暂时还没有任何复合映射。从 $\A$ 到自身的\emph{形式微分同胚} $\Phi$，是任意一列形如\eref{1.5}的映射，唯一的条件是每个 $\Phi^1$ 都是 $\hom_{\A}(X_0,X_1)$ 的线性自同构。公式\eref{1.7}使这些形式微分同胚组成一个群。现在设映射 $\mu^{\bullet}_{\A}$ 使 $\A$ 成为非含幺 $\Ai$ 范畴。给定 $\Phi$，可以递归求解\eref{1.6}，证明存在唯一的另一组 $\Ai$ 结构 $\mu^{\bullet}_{\widetilde{\A}}$（其基础向量空间相同），使 $\Phi:\A\to\widetilde{\A}$ 成为非含幺 $\Ai$ 函子。这就在所有可能的 $\mu^{\bullet}_{\A}$ 组成的集合上定义了形式微分同胚群的作用；我们记 $\widetilde{\A}=\Phi_*\A$。
\end{definition*}

\begin{note}\label{st:1.3}
“形式微分同胚”这一名称来自非交换几何；在那里，可以把 $\Ai$ 结构看作分次非交换仿射空间上次数为 $-1$ 的向量场，而上述作用就是坐标变换。
\end{note}

\subsection{自然变换}\label{sec:1d}
\begin{definition*}
非含幺 $\Ai$ 函子 $\F:\A\to\B$ 本身构成一个非含幺 $\Ai$ 范畴 $\Q=\nufun(\A,\B)$ 的对象。该范畴态射链空间中的元素 $T\in\hom^g_{\Q}(\F_0,\F_1)$，是一列 $(T^0,T^1,\ldots)$，其中每个 $T^d$ 都是对所有 $(X_0,\ldots,X_d)$ 给出的一族多重线性映射
\[
\hom_{\A}(X_{d-1},X_d)\otimes\cdots\otimes\hom_{\A}(X_0,X_1)\longrightarrow\hom_{\B}(\F_0X_0,\F_1X_d)[g-d].
\]
特别地，$T^0$ 是一族 $\hom^g_{\B}(\F_0X,\F_1X)$ 中的元素，每个 $\A$ 的对象 $X$ 对应一个。沿用 Fukaya 的术语，称这样的 $T$ 为从 $\F_0$ 到 $\F_1$ 的、次数为 $g=|T|$ 的\emph{预自然变换}。边界算子为
\end{definition*}

\begin{equation}\label{eq:1.8}
\begin{aligned}
&\mu^1_{\Q}(T)^d(a_d,\ldots,a_1)\\
&=\sum_{r,i}\sum_{s_1,\ldots,s_r}(-1)^{\dagger}\mu^r_{\B}\bigl(
\F_1^{s_r}(a_d,\ldots,a_{d-s_r+1}),\ldots,\F_1^{s_{i+1}}(\ldots,a_{s_1+\cdots+s_i+1}),\\
&\hspace{4em}T^{s_i}(a_{s_1+\cdots+s_i},\ldots,a_{s_1+\cdots+s_{i-1}+1}),
\F_0^{s_{i-1}}(a_{s_1+\cdots+s_{i-1}},\ldots),\ldots,\F_0^{s_1}(a_{s_1},\ldots,a_1)\bigr)\\
&\quad-\sum_{m,n}(-1)^{\sgnmark_n+|T|-1}T^{d-m+1}\bigl(a_d,\ldots,a_{n+m+1},
\mu^m_{\A}(a_{n+m},\ldots,a_{n+1}),a_n,\ldots,a_1\bigr).
\end{aligned}
\end{equation}
第一重求和取遍 $1\le i\le r$ 以及分拆 $s_1+\cdots+s_r=d$，其中 $s_i$ 可以为零；并且
\[
\dagger=(|T|-1)(|a_1|+\cdots+|a_{s_1+\cdots+s_{i-1}}|-s_1-\cdots-s_{i-1}).
\]
\SourcePage{19}
\begin{definition*}
上闭链 $\mu^1_{\Q}(T)=0$ 才是真正的自然变换，而上边界给出了它们之间的链同伦概念。
\end{definition*}

对于 $T_1\in\hom_{\Q}(\F_0,\F_1)$、$T_2\in\hom_{\Q}(\F_1,\F_2)$，复合公式要简单得多：
\begin{equation}\label{eq:1.9}
\begin{aligned}
&\mu^2_{\Q}(T_2,T_1)^d(a_d,\ldots,a_1)\\
&=\sum_{r,i,j}\sum_{s_1,\ldots,s_r}(-1)^{\circ}\mu^r_{\B}\bigl(
\F_2^{s_r}(a_d,\ldots,a_{d-s_r+1}),\ldots,\F_2^{s_{j+1}}(\ldots),\\
&\qquad T_2^{s_j}(a_{s_1+\cdots+s_j},\ldots,a_{s_1+\cdots+s_{j-1}+1}),
\F_1^{s_{j-1}}(\ldots),\ldots,\F_1^{s_{i+1}}(\ldots),\\
&\qquad T_1^{s_i}(a_{s_1+\cdots+s_i},\ldots,a_{s_1+\cdots+s_{i-1}+1}),
\F_0^{s_{i-1}}(\ldots),\ldots,\F_0^{s_1}(a_{s_1},\ldots,a_1)\bigr),
\end{aligned}
\end{equation}
其中
\[
\circ=\sum_{k=1}^{s_1+\cdots+s_{j-1}}(|T_2|-1)(|a_k|-1)
+\sum_{k=1}^{s_1+\cdots+s_{i-1}}(|T_1|-1)(|a_k|-1).
\]
注意，这里不再出现 $\mu^{\bullet}_{\A}$。$d>2$ 时的 $\mu^d_{\Q}$ 的公式，完全遵循\eref{1.9}的同一模式。\znote{11}{两个符号 $\dagger$ 与 $\circ$ 都是纯记账：在\eref{1.8}里要把 $T$ 所在的块移到最右端，须跨过其右侧共 $s_1+\cdots+s_{i-1}$ 个输入，每个输入贡献因子 $(-1)^{(|T|-1)(|a_k|-1)}$，累乘即得 $\dagger$；\eref{1.9}要先后把 $T_2$ 与 $T_1$ 各移一次，于是符号是两段求和。看清这一点，就不必每次重算。[译者按（上下文推断，无外部依据）：由 $T$ 的块次数 $|T|-1$ 与各输入的次数平移逐项相消得出。]}

给定 $\Ai$ 函子 $\F_0,\F_1:\A\to\B$ 之间的自然变换 $T$，考虑元素 $[T_X^0]\in\Hom_{H(\B)}(\F_0X,\F_1X)$。对于所有 $[a]\in\Hom_{H(\A)}(X_0,X_1)$，它们满足自然性条件
\[
[T^0_{X_1}]\cdot[\F_0^1(a)]=(-1)^{|a|\cdot|T|}[\F_1^1(a)]\cdot[T^0_{X_0}].
\]
因此，它们组成上同调函子 $H(\F_0),H(\F_1)$ 之间的自然变换，记作 $H(T)$。更形式化地，用
\[
\operatorname{Nu\text{-}Fun}(H(\A),H(\B))
\]
表示由非含幺线性分次函子 $H(\A)\to H(\B)$ 及其自然变换组成的范畴。那么，赋值 $\F\mapsto H(\F)$、$T\mapsto H(T)$ 定义了非含幺函子
\begin{equation}\label{eq:1.10}
H(\nufun(\A,\B))\longrightarrow\operatorname{Nu\text{-}Fun}(H(\A),H(\B)).
\end{equation}

\clearpage
\subsection{复合函子}\label{sec:1e}
与固定的 $\G:\A\to\B$ 作左复合和右复合，对任意非含幺 $\Ai$ 范畴 $\C$，定义了两个非含幺 $\Ai$ 函子
\[
\begin{aligned}
\LG_{\G}&:\nufun(\C,\A)\longrightarrow\nufun(\C,\B),\\
\RG_{\G}&:\nufun(\B,\C)\longrightarrow\nufun(\A,\C).
\end{aligned}
\]
它们在对象上的作用已作过说明。在预自然变换上的一阶映射，对于 $T\in\hom_{\nufun(\C,\A)}(\F_0,\F_1)$ 为
\begin{equation}\label{eq:1.11}
\begin{aligned}
&(\LG^1_{\G}T)^d(a_d,\ldots,a_1)\\
&=\sum_{r,i}\sum_{s_1,\ldots,s_r}(-1)^{\dagger}\G^r\bigl(
\F_1^{s_r}(a_d,\ldots,a_{d-s_r+1}),\ldots,\F_1^{s_{i+1}}(\ldots),\\
&\qquad T^{s_i}(a_{s_1+\cdots+s_i},\ldots,a_{s_1+\cdots+s_{i-1}+1}),
\F_0^{s_{i-1}}(\ldots),\ldots,\F_0^{s_1}(a_{s_1},\ldots,a_1)\bigr),
\end{aligned}
\end{equation}
\SourcePage{20}
其中 $\dagger$ 与\eref{1.8}相同；而对于 $T\in\hom_{\nufun(\B,\C)}(\F_0,\F_1)$，有
\[
\begin{aligned}
(\RG^1_{\G}T)^d(a_d,\ldots,a_1)
=\sum_r\sum_{s_1,\ldots,s_r}T^r\bigl(\G^{s_r}(a_d,\ldots,a_{d-s_r+1}),\ldots,\G^{s_1}(a_{s_1},\ldots,a_1)\bigr).
\end{aligned}
\]
高阶项 $\RG^d_{\G}$，$d\ge2$，全部为零。$\LG_{\G}$ 则并非如此，但其相应公式是\eref{1.11}的直接推广，例如见\bcite{55}，公式 (8.42)。特别地，其中不出现 $\Ai$ 范畴 $\A,\B,\C$ 的结构映射 $\mu$。复合函子本身也可复合，其行为与预期一致：
\begin{equation}\label{eq:1.12}
\RG_{\G_1}\circ\RG_{\G_2}=\RG_{\G_2\circ\G_1},\qquad
\LG_{\G_1}\circ\LG_{\G_2}=\LG_{\G_1\circ\G_2},\qquad
\LG_{\G_1}\circ\RG_{\G_2}=\RG_{\G_2}\circ\LG_{\G_1}.
\end{equation}
下面还有一个稍不显然的关系；它推广了 Gerstenhaber 关于 Hochschild 上同调上的 Yoneda 积具有分次交换性的观察。

\begin{lemma}\label{st:1.4}
设 $\Q=\nufun(\A,\A)$。假设有函子 $\G_1,\G_2\in\Ob\Q$ 和自然变换 $T_1\in\hom_{\Q}(\G_1,\Id_{\A})$、$T_2\in\hom_{\Q}(\G_2,\Id_{\A})$。那么
\[
[\mu^2_{\Q}(T_1,\LG^1_{\G_1}(T_2))]
=(-1)^{\ddagger}[\mu^2_{\Q}(T_2,\RG^1_{\G_2}(T_1))]
\in H(\hom_{\Q}(\G_1\circ\G_2,\Id_{\A})),
\]
其中 $\ddagger=(|T_1|-1)(|T_2|-1)+1$。
\end{lemma}
\begin{proof}
定义 $U\in\hom^{|T_1|+|T_2|-1}_{\Q}(\G_1\circ\G_2,\Id_{\A})$ 为
\[
\begin{aligned}
&U^d(a_d,\ldots,a_1)\\
&=\sum_{m,n,r}\sum_{s_1,\ldots,s_r}(-1)^{(|T_2|-1)\sgnmark_n}
T_1^{r+d-n-m+1}\bigl(a_d,\ldots,a_{n+m+1},\\
&\qquad T_2^m(a_{n+m},\ldots,a_{n+1}),
\G_2^{s_r}(a_n,\ldots,a_{n-s_r+1}),\ldots,\G_2^{s_1}(a_{s_1},\ldots,a_1)\bigr).
\end{aligned}
\]
直接计算可知，$\mu^1_{\Q}(U)$ 就是本引理所陈述等式两侧的差。
\end{proof}

\begin{note}\label{st:1.5}
与固定函子作复合，当然只是完整情形的一部分。为求简明，我们就讨论到这里，尽管这会使后面讨论拟等价时略有不便。
\end{note}

\clearpage
\subsection{长度滤过}\label{sec:1f}
\begin{definition*}
设 $\G_0,\G_1$ 是 $\Q=\nufun(\A,\B)$ 的对象。\emph{长度滤过}是 $\hom_{\Q}(\G_0,\G_1)$ 的一个完备递减滤过 $F^{\bullet}$：其 $F^r$ 项由满足 $T^0=\cdots=T^{r-1}=0$ 的预自然变换 $T$ 组成。
\end{definition*}

\SourcePage{21}
令 $A=H(\A)$、$B=H(\B)$、$G_k=H(\G_k)$，则相应谱序列的起始项为
\begin{equation}\label{eq:1.13}
E_1^{rs}=\prod_{X_0,\ldots,X_r}\Hom^s_{\K}\bigl(\Hom_A(X_{r-1},X_r)\otimes\cdots\otimes\Hom_A(X_0,X_1),\Hom_B(G_0X_0,G_1X_r)\bigr).
\end{equation}
即使该谱序列一般未必收敛，仍可将其用于比较论证。在最初等的层次上，由于\eref{1.13}只涉及上同调数据，有以下结论。

\begin{lemma}\label{st:1.6}
取 $\G_0,\G_1,\G_2\in\Ob\Q$，并设 $T$ 是从 $\G_0$ 到 $\G_1$ 的自然变换。假设对每个 $X\in\Ob\A$，与 $[T^0]$ 作右复合给出同构 $\Hom_{H(\B)}(\G_1X,\G_2X)\to\Hom_{H(\B)}(\G_0X,\G_2X)$。那么，在 $H(\Q)$ 中与 $T$ 作右复合给出同构 $\Hom_{H(\Q)}(\G_1,\G_2)\to\Hom_{H(\Q)}(\G_0,\G_2)$。（左复合也有相应结论。）
\end{lemma}

\begin{lemma}\label{st:1.7}
若 $\G:\A\to\B$ 是上同调全忠实的非含幺 $\Ai$ 函子，那么对任意 $\C$，$\LG_{\G}:\nufun(\C,\A)\to\nufun(\C,\B)$ 也具有这一性质。（对于 $\RG_{\G}:\nufun(\B,\C)\to\nufun(\A,\C)$ 的相应陈述一般不成立，但例如在 $\G$ 为拟同构时成立。）
\end{lemma}

由\eref{1.8}可见，\eref{1.13}上的微分 $\partial=\partial^{rs}$ 为
\[
\begin{aligned}
(\partial T)(a_{r+1},\ldots,a_1)
&=(-1)^{(r+s)(|a_1|-1)+1}T(a_{r+1},\ldots,a_2)G_0(a_1)\\
&\quad+(-1)^{\sgnmark_r+(r+s)}G_1(a_{r+1})T(a_r,\ldots,a_1)\\
&\quad+\sum_n(-1)^{\sgnmark_n+1+(r+s-1)}T(a_{r+1},\ldots,a_{n+2}a_{n+1},\ldots,a_1).
\end{aligned}
\]
（这里的约定与通常使用的约定略有不同：$T$ 是 $E_1^{rs}$ 的元素，因此作用在上同调层次上；同样，$a_i$ 是 $A$ 中的态射。）

\begin{definition*}
我们把这一双分次复形记作 $CC^{r+s}(A,B)^s=E_1^{rs}$，把它的 $\partial$ 上同调记作 $HH^{r+s}(A,B)^s=E_2^{rs}$，并分别称为 $A$ 以 $B$ 为系数的\emph{Hochschild 复形}和\emph{Hochschild 上同调}（它们依赖于函子 $G_0,G_1$，尽管记号并未显示这一点）。第一个非平凡列为
\[
E_2^{0s}=HH^s(A,B)^s=\Hom^s_{\operatorname{Nu\text{-}Fun}(A,B)}(G_0,G_1),
\]
而边缘同态 $H^s(\hom_{\Q}(\G_0,\G_1))\to E_2^{0s}$ 为 $T\mapsto H(T)$。因此，其他各列的 Hochschild 上同调群可以视为\eref{1.10}全忠实的障碍。
\end{definition*}

\SourcePage{22}
\begin{note}\label{st:1.8}
Hochschild 上同调的经典定义针对一个 $\K$ 代数 $A$ 及一个双模 $M$，给出分次群 $HH(A,M)$。这一概念可直接推广到分次代数与分次双模；此时 Hochschild 上同调成为双分次的。作为一个特殊情形，设有两个代数同态 $G_0,G_1:A\to B$。那么 $B$ 成为 $A$ 双模（左侧通过 $G_1$ 作用，右侧通过 $G_0$ 作用），因而得到 Hochschild 上同调群 $HH(A,B)$。上面考察的群，就是把这一概念从代数（可以视为只有一个对象的范畴）直接推广到一般范畴的结果，只是我们对双分次分量的标记略有些不合常规（下面将在\stref{笔记}{2.7}中，从更抽象的同调代数观点继续这一讨论）。

为避免术语混淆，要指出，在 $\Ai$ 世界中也有 Hochschild 上同调理论，即与一个 $\Ai$ 代数及其上的一个 $\Ai$ 双模相联系的分次群 $HH(\A,\M)$。同样地，如果 $\G_0,\G_1:\A\to\B$ 是 $\Ai$ 代数同态，可以借助它们使 $\B$ 成为 $\A$ 上的 $\Ai$ 双模，从而定义 $HH(\A,\B)$。把这一概念推广到 $\Ai$ 范畴，所得恰好就是我们的谱序列所趋向的上同调 $H(\hom_{\Q}(\G_0,\G_1))$。
\end{note}

\subsection{构造函子}\label{sec:1g}
接下来几段讨论 $\Ai$ 范畴与 $\Ai$ 函子的构造，它们属于同伦代数，而非同调代数。作为第一个问题，设有非含幺 $\Ai$ 范畴 $\A,\B$，其基础上同调范畴为 $A,B$，并给定一个非含幺函子 $G:A\to B$。考虑与函子对 $(G,G)$ 相联系的 Hochschild 上同调 $HH(A,B)$。

\begin{lemma}\label{st:1.9}
若对所有 $r\ge3$ 都有 $HH^2(A,B)^{2-r}=0$，那么 $G$ 可以由一个非含幺 $\Ai$ 函子 $\G:\A\to\B$ 实现，即 $H(\G)=G$。
\end{lemma}
\begin{proof}
这是一个障碍理论论证，可看作谱序列的非线性类似物。设对某个 $r\ge3$，已经定义 $\G^1,\ldots,\G^{r-1}$，使 $\Ai$ 函子方程\eref{1.6}在不超过 $r-1$ 阶时成立（当 $r=3$ 时，由于 $G$ 是函子，总能做到这一点）。任意暂取一个 $\G^r$。由此得到的 $r$ 阶 $\Ai$ 函子方程两侧之差，是一族映射
\begin{equation}\label{eq:1.14}
\theta:\hom_{\A}(X_{r-1},X_r)\otimes\cdots\otimes\hom_{\A}(X_0,X_1)\longrightarrow\hom_{\B}(\G(X_0),\G(X_r))[2-r].
\end{equation}
直接计算表明，$\theta$ 在上同调上诱导一个映射 $T$；计算从略。此外，$T$ 是 $CC^2(A,B)^{2-r}$ 中的上闭链。由假设，它在 Hochschild 复形中可写成 $\partial S$。选择一族链层次的映射 $\sigma$，使其类似于\eref{1.14}，并代表 $S$；
\SourcePage{23}
然后把已给定的 $\G^{r-1}$ 改为 $\G^{r-1}-\sigma$。这不影响 $r-1$ 阶的 $\Ai$ 函子方程，却将 $T$ 改为 $T-\partial S$，而后者按假设为零。因此，$r$ 阶 $\Ai$ 函子方程现在在上同调层次上成立。适当选择新的 $\G^r$，便可使它在链层次上也成立，这就完成了归纳步骤。
\end{proof}

我们还需要一个相对版本。
\begin{lemma}\label{st:1.10}
设 $\A,\B$ 是非含幺 $\Ai$ 范畴，$\widetilde{\A}\subset\A$ 是全子范畴，相应的上同调范畴为 $A,B,\widetilde A$。设 $G:A\to B$ 是通常的非含幺函子，$\widetilde G$ 是它在 $\widetilde A$ 上的限制。假设由此得到的 Hochschild 上同调群之间的自然限制映射
\[
HH^{r+s}(A,B)^s\longrightarrow HH^{r+s}(\widetilde A,B)^s
\]
当 $r+s=2$、$r\ge3$ 时是单射，当 $r+s=1$、$r\ge2$ 时是满射。那么，给定任意满足 $H(\widetilde{\G})=\widetilde G$ 的非含幺 $\Ai$ 函子 $\widetilde{\G}:\widetilde{\A}\to\B$，都可以将其扩张为满足 $H(\G)=G$ 的 $\G:\A\to\B$。
\end{lemma}
\begin{proof}
为此，使用由 $CC(A,B)$ 中限制到 $CC(\widetilde A,B)$ 后为零的链组成的子复形。假设保证这一子复形在 $r\ge3$ 的双次数 $(2,2-r)$ 上无上同调。作此修改后，证明与前面完全相同。
\end{proof}

\clearpage
\subsection{同伦}\label{sec:1h}
\begin{definition*}
考虑 $\Q=\nufun(\A,\B)$ 的对象集中的两个非含幺 $\Ai$ 函子 $\F_0,\F_1$，假设它们在对象上的作用相同。令 $D=\F_0-\F_1\in\hom^1_{\Q}(\F_0,\F_1)$ 为由 $D^0=0$、$D^d=\F_0^d-\F_1^d$（$d>0$）定义的预自然变换。它是真正的自然变换，即 $\mu^1_{\Q}(D)=0$。如果存在还满足 $T^0=0$ 的 $T\in\hom^0_{\Q}(\F_0,\F_1)$，使 $D=\mu^1_{\Q}(T)$，就称 $\F_0,\F_1$ \emph{同伦}。至少从表面上看，这一名称是合理的，因为同伦函子在上同调上诱导相同的函子：$H(\F_0)=H(\F_1)$。\znote{15}{要求 $T^0=0$ 并非技术洁癖：同伦的差 $D=\F_0-\F_1$ 在零阶已为零，若允许 $T^0\neq0$，则 $\mu^1_{\Q}(T)$ 的零阶项会混入 $\F_0^1-\F_1^1$ 的修正量，$H(\F_0)=H(\F_1)$ 便不再自动成立。把 $T^0$ 归零后，$\mu^1_{\Q}(D)=0$ 逐阶给出 $\F_0^1-\F_1^1$ 是上边界，两个函子在上同调上因而一致。[来源：Seidel, Fukaya categories and Picard–Lefschetz theory, §2c（同伦与同构）；并见\stref{笔记}{1.11}。]}与固定函子作左复合、右复合，都保持同伦。这意味着，若 $T$ 是从 $\F_0$ 到 $\F_1$ 的同伦，那么 $\RG^1_{\G}(T)$ 是 $\F_0\circ\G$ 与 $\F_1\circ\G$ 之间的同伦，而 $\LG^1_{\G}(T)$ 是 $\G\circ\F_0$ 与 $\G\circ\F_1$ 之间的同伦。
\end{definition*}

同伦是等价关系。原则上这不难验证，但计算可能显得有些混乱，原因如下：由于考察的函子 $\F_k$ 在对象上的作用相同，对于所有 $k,l$，分次向量空间 $\hom_{\Q}(\F_k,\F_l)$ 都相同，而我们需要把预自然变换从一个空间移到另一个空间中。为避免混淆，记 $T_{[kl]}$ 表示把 $T$ 视为 $\hom_{\Q}(\F_k,\F_l)$ 中的元素。采用这一记号，基本观察是：对于任意三个函子 $\F_j,\F_k,\F_l$，
\begin{equation}\label{eq:1.15}
\begin{aligned}
\mu^1_{\Q}(T_{[jl]})
&=\mu^1_{\Q}(T_{[jk]})-\mu^2_{\Q}(\F_k-\F_l,T_{[jk]})\\
&=\mu^1_{\Q}(T_{[kl]})+\mu^2_{\Q}(T_{[kl]},\F_j-\F_k).
\end{aligned}
\end{equation}
\SourcePage{24}
为证明传递性，设 $T_1$ 是从 $\F_0$ 到 $\F_1$ 的同伦，$T_2$ 是从 $\F_1$ 到 $\F_2$ 的同伦。由\eref{1.15}可见
\begin{equation}\label{eq:1.16}
T=(T_1)_{[02]}+(T_2)_{[02]}+\mu^2_{\Q}(T_2,T_1)
\end{equation}
满足
\[
\begin{aligned}
\mu^1_{\Q}(T)
&=\F_1-\F_2-\mu^2_{\Q}(\F_1-\F_2,T_1)\\
&\quad+\F_0-\F_1+\mu^2_{\Q}(T_2,\F_0-\F_1)+\mu^1_{\Q}\mu^2_{\Q}(T_2,T_1)
=\F_0-\F_2,
\end{aligned}
\]
因此给出从 $\F_0$ 到 $\F_2$ 的同伦。对称性的证明也类似：设 $T_1$ 是从 $\F_0$ 到 $\F_1$ 的同伦，考虑映射
\begin{equation}\label{eq:1.17}
\begin{aligned}
\phi &: \hom_{\Q}(\F_1,\F_0)\longrightarrow\hom_{\Q}(\F_1,\F_1),
&\phi(T)&=(-1)^{|T|}\mu^2_{\Q}(T_1,T)+T_{[11]},\\
\psi &: \hom_{\Q}(\F_1,\F_1)\longrightarrow\hom_{\Q}(\F_0,\F_1),
&\psi(T)&=\mu^2_{\Q}(T,T_1)+T_{[01]}.
\end{aligned}
\end{equation}
由于各式中的最后一项，它们是同构；再应用\eref{1.15}，可见它们还是链映射。最后，再次利用\eref{1.15}，得到
\[
\begin{aligned}
\phi((\F_0-\F_1)_{[10]})&=\mu^1_{\Q}((T_1)_{[11]}),\\
\psi(\mu^1_{\Q}((T_1)_{[11]}))&=\mu^1_{\Q}\mu^2_{\Q}((T_1)_{[11]},T_1)+(\F_0-\F_1).
\end{aligned}
\]
由此可知，$-T_1-\mu^2_{\Q}((T_1)_{[11]},T_1)$ 在 $\psi\circ\phi$ 下的原像，是从 $\F_1$ 到 $\F_0$ 的链同伦。

\begin{note}\label{st:1.11}
当 $\B$ 具有严格恒等态射时，上述计算及一般的同伦概念会变得更清楚一些；见后面的\secref{2c}。这里先给出一个来自有理同伦论的不同解释，参看\bcite{68}，第 X 章。考虑如下 dg 代数：$I=\K u_0\oplus\K u_1\oplus\K h$，其中 $|u_0|=|u_1|=0$、$|h|=1$；微分和乘法为（省略因次数原因而为零的式子）$\partial u_0=h=-\partial u_1$，以及
\[
\begin{gathered}
u_0u_0=u_0,\qquad u_1u_1=u_1,\qquad u_0u_1=u_1u_0=0,\\
u_1h=h,\qquad hu_1=u_0h=0,\qquad hu_0=h.
\end{gathered}
\]
$H(I)\cong\K$，生成元为 $u_0+u_1$，故投影到 $u_k$ 所生成子空间的两个映射 $f_k:I\to\K$（$k=0,1$）都是拟同构。若 $\B$ 是非含幺 $\Ai$ 代数，很容易定义张量积 $I\otimes\B$：在基础分次向量空间的张量积上，令
\[
\mu^1(i\otimes b)=i\otimes\mu^1_{\B}(b)+(-1)^{|i|+|b|}\partial(i)\otimes b,
\]
而对于 $d\ge2$，令
\[
\mu^d(i_d\otimes b_d,\ldots,i_1\otimes b_1)
=(-1)^{\triangleleft}i_d\cdots i_1\otimes\mu^d_{\B}(b_d,\ldots,b_1),
\qquad\triangleleft=\sum_{j>k}(|b_j|-1)|i_k|.
\]
简单计算表明，$T$ 是两个 $\Ai$ 同态 $\F_0,\F_1:\A\to\B$ 之间的同伦，当且仅当
\[
\begin{aligned}
\G(a_d,\ldots,a_1)
&=u_0\otimes\F_0(a_d,\ldots,a_1)+u_1\otimes\F_1(a_d,\ldots,a_1)\\
&\quad+(-1)^{\sgnmark_d}h\otimes T(a_d,\ldots,a_1)
\end{aligned}
\]
\SourcePage{25}
定义了一个 $\Ai$ 同态 $\G:\A\to I\otimes\B$。这一情形可归结为下面的交换图。
\SourceDiagram{diagram-1h}{0.24715358}
作出显然的修改后，同样的叙述也适用于 $\Ai$ 范畴与函子。
\end{note}

\subsection{同调扰动理论}\label{sec:1i}
设 $\B$ 是非含幺 $\Ai$ 范畴。假设对每对对象 $(X_0,X_1)$，有如下图：
\SourceDiagram{diagram-1i}{0.40292265}
其中 $(\hom_{\A}(X_0,X_1),\mu^1_{\A})$ 是某个向量空间链复形；$\F^1,\G^1$ 是链映射；$T^1$ 的次数为 $-1$，且满足 $\mu^1_{\B}T^1+T^1\mu^1_{\B}=\F^1\G^1-\id$。\emph{扰动引理}即以下陈述。\znote{17}{这正是经典的同伦扰动引理在 $\Ai$ 范畴上的形式：若 $f:X\to Y$、$\nabla:Y\to X$ 为链映射，$\Phi$ 为满足 $f\nabla=\id$、$\nabla f=\id+d\Phi+\Phi d$ 的收缩同伦，则扰动后的微分可由显式公式 $\Phi^{\tau}=\sum_{r\ge0}\Phi(1+d^{\tau}\Phi)^{r}$ 给出。原书的 $\F,\G,T$ 分别对应 $f,\nabla,\Phi$，\eref{1.18}即该显式公式的展开。[来源：nLab \textquotedblleft homological perturbation theory\textquotedblright，https://\allowbreak ncatlab.org/\allowbreak nlab/\allowbreak show/\allowbreak homological+\allowbreak perturbation+\allowbreak theory（2026-11-15 访问）。]}

\begin{proposition}\label{st:1.12}
可以构造一个非含幺 $\Ai$ 范畴 $\A$，其对象集为 $\Ob\A=\Ob\B$，态射空间为 $\hom_{\A}(X_0,X_1)$，一阶结构映射为 $\mu^1_{\A}$；同时得到非含幺 $\Ai$ 函子 $\F:\A\to\B$、$\G:\B\to\A$，它们在对象上均为恒等，一阶部分是给定的映射 $\F^1,\G^1$；还得到 $\F\circ\G$ 与 $\Id_{\B}$ 之间的同伦 $T$，其起始项是给定的 $T^1$。
\end{proposition}

所有这些结构都由明确公式定义，通常写成递归形式。例如，
\begin{equation}\label{eq:1.18}
\begin{aligned}
\F^d(a_d,\ldots,a_1)
&=\sum_r\sum_{s_1,\ldots,s_r}T^1\bigl(\mu^r_{\B}(\F^{s_r}(a_d,\ldots,a_{d-s_r+1}),\ldots,\F^{s_1}(a_{s_1},\ldots,a_1))\bigr),\\
\mu^d_{\A}(a_d,\ldots,a_1)
&=\sum_r\sum_{s_1,\ldots,s_r}\G^1\bigl(\mu^r_{\B}(\F^{s_r}(a_d,\ldots,a_{d-s_r+1}),\ldots,\F^{s_1}(a_{s_1},\ldots,a_1))\bigr),
\end{aligned}
\end{equation}
其中两个和式都对 $r\ge2$ 的分拆 $s_1+\cdots+s_r=d$ 求和。很容易验证，$\mu^{\bullet}_{\A}$ 确实是 $\Ai$ 结构，$\F$ 确实是 $\Ai$ 函子。\znote{17}{把 $\hom_{\A}$ 取成 $\hom_{\B}$ 的上同调（\stref{命题}{1.12}中\textquotedblleft 分裂\textquotedblright 的极端情形），即得 Kadeishvili 定理：dg 代数的上同调自然带有 $\Ai$ 代数结构，且与原代数弱等价。这一转移定理正是镜像对称中\textquotedblleft Floer 上同调仍带 $\Ai$ 结构\textquotedblright 的来源，也是本书后面以上同调范畴为出发点的依据。[来源：nLab \textquotedblleft A-infinity-algebra\textquotedblright（Kadeishvili–Merkulov 定理一节），https:/\allowbreak /\allowbreak ncatlab.org/\allowbreak nlab/\allowbreak show/\allowbreak A-infinity-algebra（2026-11-15 访问）。]}
\SourcePage{26}
为了使计算更具几何意味，可以把递归定义“展开”，写成涉及树的求和的直接表达式（更精确地，是\secref{9d}意义下的稳定树）。$\G$ 和 $T$ 的构造也可以这样处理，但相应公式略为复杂一些。

\begin{note}\label{st:1.13}
一个典型应用如下。给定非含幺 $\Ai$ 范畴 $\B$，把每个复形 $(\hom_{\B}(X_0,X_1),\mu^1_{\B})$ 分裂为微分为零的直和项与一个无上同调的补项。取 $\hom_{\A}(X_0,X_1)$ 为前一个直和项，$\F^1$ 为包含映射，$\G^1$ 为投影。为无上同调的补项选择一个收缩同伦，就得到所需的 $T^1$。于是\stref{命题}{1.12}构造出满足 $\mu^1_{\A}=0$、且与 $\B$ 拟同构的 $\A$。在这一特殊情形中，$\G\circ\F$ 也同伦于 $\Id_{\A}$。为说明这一点，先注意，由于 $\F\circ\G$ 同伦于 $\Id_{\B}$，$\G\circ\F\circ\G$ 同伦于 $\G$。\stref{引理}{1.7}说明
\[
\RG^1_{\G}:\hom_{\nufun(\A,\A)}(\Id_{\A},\G\circ\F)
\longrightarrow\hom_{\nufun(\B,\A)}(\G,\G\circ\F\circ\G)
\]
是链复形的拟同构。此外，改为考察零阶项为零的预自然变换组成的子复形，这一结论仍然成立。显然，$\RG^1_{\G}$ 将 $\Id_{\A}-\G\circ\F$ 映为 $\G-\G\circ\F\circ\G$，因此这就证明了所需的第二个同伦的存在性。
\end{note}

\begin{corollary}\label{st:1.14}
非含幺 $\Ai$ 范畴之间的任意拟同构，都具有同伦意义下的逆。
\end{corollary}
\begin{proof}
设给定的拟同构为 $\mathcal K:\B\to\widetilde{\B}$。应用扰动引理，构造由非含幺 $\Ai$ 范畴及拟同构组成的如下图：
\SourceDiagram{diagram-1-14}{0.22638437}
其中 $\mu^1_{\A}=\mu^1_{\widetilde{\A}}=0$。复合 $\mathcal H=\widetilde{\G}\circ\mathcal K\circ\F$ 是两个一阶复合为零的 $\Ai$ 范畴之间的拟同构，因此必然是形式微分同胚，从而相对于复合具有严格的逆 $\mathcal H^{-1}$。考虑 $\mathcal I=\F\circ\mathcal H^{-1}\circ\widetilde{\G}$。按构造，$\mathcal K$ 同伦于 $\widetilde{\F}\circ\mathcal H\circ\G=\widetilde{\F}\circ\widetilde{\G}\circ\mathcal K\circ\F\circ\G$。因此 $\mathcal K\circ\mathcal I$ 同伦于 $\widetilde{\F}\circ\mathcal H\circ\G\circ\F\circ\mathcal H^{-1}\circ\widetilde{\G}$，而后者又同伦于恒等函子。另一个同伦以同样方式构造。
\end{proof}

\SourcePage{27}
\begin{note}\label{st:1.15}
即使不计拓扑学中的先驱结果\bcite{137}，\stref{命题}{1.12}也有很长的历史。早期特殊情形见\bcite{80}和\bcite{69}；经过此后种种推广与简化，最终得到这里陈述的版本，取自\bcite{98}。证明、更多参考文献，以及一幅鲤鱼图，都可参看该文。至少在 $\operatorname{char}(\K)=0$ 时，可以由控制 $\Ai$ 结构的操作子的“余纤维性”推出抽象扰动引理的存在，见\bcite{99}。\stref{推论}{1.14}最早的出处似乎是\bcite{81}、\bcite{111}。对于 $L_\infty$ 代数，略弱一些的类似结果出现于\bcite{88}，第 4.5 节；其证明与上面基本相同。若想了解体现\stref{引理}{1.9}思想的障碍理论方法，可参看\bcite{60}，附录 A5。
\end{note}

\clearpage
\subsection{\texorpdfstring{$A_\infty$}{A∞} 模}\label{sec:1j}
\begin{definition*}
设 $\Ch$ 是 $\K$ 向量空间链复形组成的 dg 范畴，并将其视为 $\Ai$ 范畴。对于固定的非含幺 $\Ai$ 范畴 $\A$，将 $\A$ 上非含幺（右）$\Ai$ 模组成的非含幺 $\Ai$ 范畴定义为
\[
\operatorname{nu\text{-}mod}(\A)=\nufun(\A^{\mathrm{opp}},\Ch).
\]
\end{definition*}

\begin{definition*}
有必要具体写出细节。$\Q=\operatorname{nu\text{-}mod}(\A)$ 的对象 $\M$，由对每个 $X\in\Ob\A$ 给出的分次向量空间 $\M(X)$，以及对所有 $d\ge1$ 给出的映射组成：
\[
\mu^d_{\M}:\M(X_{d-1})\otimes\hom_{\A}(X_{d-2},X_{d-1})\otimes\cdots\otimes\hom_{\A}(X_0,X_1)
\longrightarrow\M(X_0)[2-d].
\]
若从 $\Ai$ 函子的角度理解，$\mu^1_{\M}$ 就是像对象 $\M(X_0)\in\Ob\Ch$ 上的微分，而对 $d\ge2$，$\mu^d_{\M}$ 是函子在态射上的作用；更精确地，是把自变量顺序反转后的 $\M^{d-1}$。方程\eref{1.6}变为
\begin{equation}\label{eq:1.19}
\begin{aligned}
&\sum_n(-1)^{\sgnmark_n}\mu^{n+1}_{\M}\bigl(\mu^{d-n}_{\M}(b,a_{d-1},\ldots,a_{n+1}),a_n,\ldots,a_1\bigr)\\
&\quad+\sum_{m,n}(-1)^{\sgnmark_n}\mu^{d-m+1}_{\M}\bigl(b,a_{d-1},\ldots,a_{n+m+1},\mu^m_{\A}(a_{n+m},\ldots,a_{n+1}),a_n,\ldots,a_1\bigr)=0,
\end{aligned}
\end{equation}
其中第二个和式满足 $n+m<d$。\znote{19}{把右模写成 $\nufun(\A^{\mathrm{opp}},\Ch)$，是把\textquotedblleft 模 $=$ 单对象范畴上的函子\textquotedblright 这一经典对应推广到 $\Ai$ 情形：$\Ai$ 代数就是只有一个对象的 $\Ai$ 范畴，其上的模便是定义在反范畴上、取值于链复形范畴的 $\Ai$ 函子。也正因如此，\eref{1.19}中 $\M(X_{d-1})$ 排在最左、$a$ 的指标递减，与 $\hom_{\A}$ 的复合顺序正好相反。[来源：nLab \textquotedblleft A-infinity-category\textquotedblright，https:/\allowbreak /\allowbreak ncatlab.org/\allowbreak nlab/\allowbreak show/\allowbreak A-infinity-category（2026-11-15 访问）。]}
\end{definition*}

\begin{definition*}
再设 $\M_0,\M_1$ 为两个非含幺 $\Ai$ 模。$\hom_{\Q}(\M_0,\M_1)$ 中的态射称为\emph{预模同态}，写成一列映射
\[
t^d:\M_0(X_{d-1})\otimes\hom_{\A}(X_{d-2},X_{d-1})\otimes\cdots\otimes\hom_{\A}(X_0,X_1)
\longrightarrow\M_1(X_0)[|t|-d+1],\qquad d\ge1.
\]
它与预自然变换 $T$ 的一般表述之间，通过显然的自变量调换相联系，同时作一个为方便而选取的符号调整：
\begin{equation}\label{eq:1.20}
t^d(b,a_{d-1},\ldots,a_1)=(-1)^{\S} T^{d-1}(a_1,\ldots,a_{d-1})(b),
\end{equation}
\SourcePage{28}
其中 $T^{d-1}(a_1,\ldots,a_{d-1})\in\hom_{\Ch}(\M_0(X_{d-1}),\M_1(X_0))$，且 $\S=(|T|-1)|b|+|T|(|T|-1)/2$。真正的模同态是满足 $\mu^1_{\Q}(t)=0$ 的那些 $t$。
\end{definition*}

为看清其具体含义，写出 $\Q$ 中的复合映射：
\begin{equation}\label{eq:1.21}
\begin{aligned}
&(\mu^1_{\Q}t)^d(b,a_{d-1},\ldots,a_1)\\
&=\sum_n(-1)^{\ddagger}\mu^{n+1}_{\M_1}\bigl(t^{d-n}(b,a_{d-1},\ldots,a_{n+1}),a_n,\ldots,a_1\bigr)\\
&\quad+\sum_n(-1)^{\ddagger} t^{n+1}\bigl(\mu^{d-n}_{\M_0}(b,a_{d-1},\ldots,a_{n+1}),a_n,\ldots,a_1\bigr)\\
&\quad+\sum_{m,n}(-1)^{\ddagger} t^{d-m+1}\bigl(b,a_{d-1},\ldots,a_{n+m+1},\mu^m_{\A}(a_{n+m},\ldots,a_{n+1}),a_n,\ldots,a_1\bigr);\\[2pt]
&(\mu^2_{\Q}(t_2,t_1))^d(b,a_{d-1},\ldots,a_1)\\
&=\sum_n(-1)^{\ddagger} t_2^{n+1}\bigl(t_1^{d-n}(b,a_{d-1},\ldots,a_{n+1}),a_n,\ldots,a_1\bigr);\\
&\mu^d_{\Q}=0\qquad\text{对所有 }d\ge3.
\end{aligned}
\end{equation}
这里 $\ddagger=|a_{n+1}|+\cdots+|a_{d-1}|+|b|-d+n+1$。$\mu^3_{\Q},\mu^4_{\Q},\ldots$ 为零，是一个非正式但有用的一般原则的实例：$\nufun(\A,\B)$ 继承的是 $\B$ 的性质，而非 $\A$ 的性质。

在上同调层次上，任意 $\M$ 都给出一个非含幺 $H(\A)$ 模 $H(\M)$；更形式化地说，它是从 $H(\A)$ 到分次向量空间范畴的非含幺反变函子。它由 $\M(X)$ 关于 $\partial(b)=(-1)^{|b|}\mu^1_{\M}(b)$ 的上同调，以及 $b\cdot a=(-1)^{|a|}\mu^2_{\M}(b,a)$ 所诱导的模结构组成。同样，给定 $\Ai$ 模同态 $t$，得到诱导映射 $H(t):H(\M_0(X))\to H(\M_1(X))$，$[b]\mapsto[(-1)^{|b|}t^1(b)]$；它们组成通常意义下的模同态。下面是\stref{引理}{1.6}的一个特殊情形。

\begin{lemma}\label{st:1.16}
设 $t:\M_0\to\M_1$ 是 $\Ai$ 模同态，并且对每个 $X\in\Ob\A$，映射 $H(t):H(\M_0(X))\to H(\M_1(X))$ 都是同构。那么，对于任意 $\N$，与 $t$ 作右复合及左复合，分别诱导拟同构
\[
\hom_{\Q}(\M_1,\N)\longrightarrow\hom_{\Q}(\M_0,\N),\qquad
\hom_{\Q}(\N,\M_0)\longrightarrow\hom_{\Q}(\N,\M_1).
\]
\end{lemma}

特别地，如果非含幺 $\Ai$ 模 $\M$ 对所有 $X$ 都满足 $H(\M(X),\mu^1_{\M})=0$，那么对所有 $\N$ 都有 $H(\hom_{\Q}(\M,\N))=H(\hom_{\Q}(\N,\M))=0$。借用\bcite{136}、\bcite{21}中的术语，可以说在 $\Ai$ 模的世界中，每个对象同时都是 K-投射的和 K-内射的。

\subsection{拉回}\label{sec:1k}
\begin{definition*}
任意非含幺 $\Ai$ 函子 $\G:\A\to\B$ 都诱导一个拉回函子
\begin{equation}\label{eq:1.22}
\G^*=\RG_{\G^{\mathrm{opp}}}:\operatorname{nu\text{-}mod}(\B)\longrightarrow\operatorname{nu\text{-}mod}(\A).
\end{equation}
\SourcePage{29}
与任意右复合一样，它的高阶项为零，因此是 dg 范畴之间的非含幺 dg 函子。具体地，
\[
\begin{aligned}
(\G^*\M)(X)&=\M(\G(X)),\\
\mu^d_{\G^*\M}(b,a_{d-1},\ldots,a_1)
&=\sum_r\sum_{s_1,\ldots,s_r}\mu^{r+1}_{\M}\bigl(b,\G^{s_r}(a_{d-1},\ldots,a_{d-s_r}),\ldots,\G^{s_1}(a_{s_1},\ldots,a_1)\bigr),\\
(\G^*t)^d(b,a_{d-1},\ldots,a_1)
&=\sum_r\sum_{s_1,\ldots,s_r}t^{r+1}\bigl(b,\G^{s_r}(a_{d-1},\ldots,a_{d-s_r}),\ldots,\G^{s_1}(a_{s_1},\ldots,a_1)\bigr).
\end{aligned}
\]
\end{definition*}

\subsection{Yoneda 嵌入}\label{sec:1l}
\begin{definition*}
对 $\A$ 的每个对象 $Y$，都有相应的非含幺 $\Ai$ 模 $\Y$，定义为 $\Y(X)=\hom_{\A}(X,Y)$、$\mu^d_{\Y}=\mu^d_{\A}$。这是一个典范的非含幺 $\Ai$ 函子 $\mathcal I=\mathcal I_{\A}:\A\to\Q=\operatorname{nu\text{-}mod}(\A)$ 的一部分，该函子称为\emph{Yoneda 嵌入}。对于 $c\in\hom_{\A}(Y_0,Y_1)$，$\mathcal I^1(c)\in\hom_{\Q}(\Y_0,\Y_1)$ 是如下预模同态：
\[
\begin{aligned}
\Y_0(X_{d-1})\otimes\hom_{\A}(X_{d-2},X_{d-1})\otimes\cdots\otimes\hom_{\A}(X_0,X_1)&\longrightarrow\Y_1(X_0),\\
(b,a_{d-1},\ldots,a_1)&\longmapsto\mu^{d+1}_{\A}(c,b,a_{d-1},\ldots,a_1).
\end{aligned}
\]
高阶项 $\mathcal I^d$，$d\ge2$，由这一公式的显然类似式定义。\znote{21}{Yoneda 嵌入给出的模 $\Y=\hom_{\A}(\cdot,Y)$ 是\textquotedblleft 自由模\textquotedblright：$\Ai$ 版 Yoneda 引理说明 $\hom_{\Q}(\Y,\M)$ 的上同调就是 $H(\M(Y))$，\eref{1.25}中的 $\lambda$ 正是这一对应的显式形式。一个常见推论是每个线性 $\Ai$ 范畴都 $\Ai$ 等价于某个 dg 范畴；但该等价只在抽象层次上成立，对 Fukaya 范畴直接使用并不现实。[来源：nLab \textquotedblleft A-infinity-category\textquotedblright，https:/\allowbreak /\allowbreak ncatlab.org/\allowbreak nlab/\allowbreak show/\allowbreak A-infinity-category（2026-11-15 访问）。]}
\end{definition*}

Yoneda 嵌入与拉回在如下意义下相容。设 $\F:\A\to\B$ 为非含幺 $\Ai$ 函子，则存在典范自然变换
\begin{equation}\label{eq:1.23}
T:\mathcal I_{\A}\longrightarrow\F^*\circ\mathcal I_{\B}\circ\F.
\end{equation}
其一阶项 $T^1$ 对每个 $Y$ 给出一个模同态，而这个模同态本身是一列映射
\begin{equation}\label{eq:1.24}
\begin{aligned}
(\mathcal I_{\A}(Y))(X_{d-1})\otimes\hom_{\A}(X_{d-2},X_{d-1})\otimes\cdots\otimes\hom_{\A}(X_0,X_1)\\
\longrightarrow (\F^*\circ\mathcal I_{\B}\circ\F(Y))(X_0).
\end{aligned}
\end{equation}
按定义，$(\mathcal I_{\A}(Y))(X_{d-1})=\hom_{\A}(X_{d-1},Y)$，且 $(\F^*\circ\mathcal I_{\B}\circ\F)(Y)(X_0)=\hom_{\B}(\F(X_0),\F(Y))$；据此，\eref{1.24}由 $(b,a_{d-1},\ldots,a_1)\mapsto(-1)^{\sgnmark_{d-1}+|b|}\F^d(b,a_{d-1},\ldots,a_1)$ 给出。类似的表达式定义 $T^r$，$r\ge2$。

从另一个方向，可以把 $\mathcal I^1$ 推广为链映射
\begin{equation}\label{eq:1.25}
\begin{aligned}
\lambda=\lambda_{\M}&:\M(Y)\longrightarrow\hom_{\Q}(\Y,\M),\\
\lambda(c)^d(b,a_{d-1},\ldots,a_1)&=\mu^{d+1}_{\M}(c,b,a_{d-1},\ldots,a_1),
\end{aligned}
\end{equation}
其中 $\M$ 是任意非含幺 $\Ai$ 模。在上同调层次上，这一映射相对于两侧的乘法具有以下容易验证的性质：
\begin{equation}\label{eq:1.26}
[\lambda_{\M}(\mu^2_{\M}(c,b))]
=[\mu^2_{\Q}(\lambda_{\M}(c),\mathcal I^1(b))]\in\Hom_{H(\Q)}(\Y_0,\M),
\end{equation}
\SourcePage{30}
其中 $[c]\in H(\M(Y_1))$、$[b]\in\Hom_{H(\A)}(Y_0,Y_1)$；以及
\begin{equation}\label{eq:1.27}
[\mu^2_{\Q}(t,\lambda_{\M_0}(c))]
=[\lambda_{\M_1}(t^1(c))]\in\Hom_{H(\Q)}(\Y,\M_1),
\end{equation}
其中 $[t]\in\Hom_{H(\Q)}(\M_0,\M_1)$、$[c]\in H(\M_0(Y))$。注意，如果只考虑来自 Yoneda 嵌入的模，并相应地令 $t$ 为 $\A$ 中某个态射的像，那么这两个等式都归结为 $\mathcal I$ 的 $\Ai$ 函子方程的特殊情形。

\begin{note}\label{st:1.17}
映射 $\lambda$ 可以视为更广泛构造的一部分，这也解释了上面列出的性质。把 $\Q$ 的 Yoneda 嵌入与沿 $\A$ 的 Yoneda 嵌入作拉回相结合，得到非含幺 dg 函子，称为\emph{线性化函子}：
\[
\LG:\Q\xrightarrow{\mathcal I_{\Q}}\operatorname{nu\text{-}mod}(\Q)
\xrightarrow{\mathcal I_{\A}^*}\Q.
\]
（这一名称来自如下事实：每个模同态在 $\LG$ 下的像，只有一阶线性项可能非零。）这个函子带有典范自然变换 $L:\Id_{\Q}\to\LG$。在一阶，它对每个 $\Ai$ 模 $\M$ 给出模同态 $l=L^1_{\M}:\M\to\M_{\mathrm{lin}}=\LG(\M)$；后者的起始项由对所有 $Y$ 给出的线性映射 $l^1:\M(Y)\to\M_{\mathrm{lin}}(Y)$ 组成。按定义，$\M_{\mathrm{lin}}(Y)=\hom_{\Q}(\Y,\M)$，其中 $\Y$ 是 $Y$ 在 Yoneda 嵌入下的像，而该线性映射与我们的 $\lambda$ 只差一个符号 $(-1)^{|c|}$。$l$ 是 $\Ai$ 模同态这一事实蕴含\eref{1.26}；$L$ 是自然变换这一事实蕴含\eref{1.27}。
\end{note}
